% ============================================================================
% sec8_robustness.tex — Paper 2: Robustness
% ============================================================================

\section{Robustness}
\label{sec:robustness}

The companion paper \citep{genicolomartins2026screening} reports
extensive reduced-form robustness checks (IQR threshold variants,
clustering alternatives, matching estimators, sensitivity to
unobservables). Here we focus on checks specific to the structural
model and the welfare counterfactuals.

\paragraph{Reduced-form stability.}
Table~\ref{tab:robustness_summary} collects the price association
across FL-definition variants, sample restrictions, and estimation
methods. The coefficient ranges from 0.036 (cross-fit) to 0.077
(CEM), significant in all specifications. The cross-fit average
breaks any mechanical link between the FL rule and the price outcome.
\citet{cinelli2020making} sensitivity analysis yields $RV_{q=1} =
17.5\%$: an unobserved confounder would need to explain 17.5\% of
residual variation in both FL presence and log prices to nullify the
association.

\begin{table}[htbp]
\centering
\caption{Robustness summary: FL--price association}
\label{tab:robustness_summary}
\small
\begin{tabular}{lccc}
\hline\hline
Check & Coeff.\ & SE & Note \\
\hline
OLS (item + year + PBU FE) & 0.064 & 0.020 & Baseline \\
Cross-fit average & 0.036 & --- & Out-of-sample \\
CEM matching & 0.077 & 0.024 & Non-parametric \\
IPW matching & 0.055 & 0.021 & Non-parametric \\
$n$-conditional (structural) & 0.062 & 0.019 & Controls for $n_{\text{genuine}}$ \\
IV (leave-one-out) & 0.194 & 0.077 & Upper bound \\
\hline
IQR 1.0$\times$ & 0.079 & 0.023 & 3,442 FL firms \\
IQR 2.0$\times$ & 0.060 & 0.022 & 2,093 FL firms \\
IQR 3.0$\times$ & 0.050 & 0.024 & 1,456 FL firms \\
\hline\hline
\end{tabular}
\end{table}

\paragraph{Structural robustness.}
Three checks address the structural estimates. \emph{First}, the
regime selection is insensitive to the estimation sample: restricting
to preg\~ao-only or convite-only subsamples yields
$\Delta$BIC~$< -10{,}000$ in favor of Regime~2 in both cases.
\emph{Second}, the dispersion ratio $\hat{\sigma}_c / \hat{\sigma}_g$
is stable across IQR multiplier thresholds: 0.70 at $1.0\times$,
0.72 at $1.5\times$ (baseline), and 0.74 at $3.0\times$---the
structural finding does not depend on the classification boundary.
\emph{Third}, the ground-truth tests
(Section~\ref{sec:results_detection}) confirm that the structural
LR score (AUC~$= 0.56$) does not improve on the simple IQR rule
(AUC~$= 0.94$), indicating that the model characterizes the
mechanism but does not substitute for the participation-based screen.

\paragraph{Random-label placebo.}
To verify that the structural findings are specific to FL firms and
not an artifact of selecting any subset of high-participation losers,
we assign the ``FL'' label to 2,735 randomly chosen non-FL firms
(matching the true FL count) and re-estimate the dispersion ratio.
Across 100 placebo draws, the mean placebo ratio is
$\hat{\sigma}_c / \hat{\sigma}_g = 1.00$ (SD~$= 0.01$)---centered
at one, as expected when labels carry no information. The true ratio
(0.72) falls below every placebo draw ($p < 0.01$). The Bajari--Ye
pairwise product under placebo labels shows no tender-FE reversal
(mean product $\approx 0.8$, compared to 0.38 for true FL firms).
The placebo indicates that the dispersion pattern and the tender-FE
reversal are specific to FL-labeled firms, not a mechanical
consequence of subsetting.

\paragraph{Latent class validation.}
An unsupervised two-component Gaussian mixture on 8,275 always-losers
(using five bid-level features, \emph{without} FL labels) discovers
a ``cover-bidder'' class with mean participation of 45.2 tenders
versus 9.0 for the ``genuine'' class. The cover class contains
71.2\% FL firms versus 20.0\% in the genuine class
(Table~\ref{tab:lca_validation}, Appendix~\ref{app:robustness}).
The unsupervised model independently recovers the defining
characteristic of the FL classification, providing structural support
for the IQR rule.

\paragraph{Staggered DiD.}
A \citet{callaway2021difference} DiD and a stacked regression
\citep{cengiz2019effect} both produce insignificant ATTs
(C\&S: 0.014, SE $= 0.039$; stacked: $-0.006$, SE $= 0.014$).
The stacked 95\% CI ($-0.034$, $0.022$) excludes the OLS baseline
(0.064), consistent with failure of parallel trends---the pre-trend
coefficients at $t = -2$ and $t = -3$ confirm that FL entry is
endogenous to market conditions. The DiD designs are invalid for
this setting, reinforcing the diagnostic framing. Full results in
Appendix~\ref{app:did}.

\paragraph{Oversight heterogeneity.}
The FL--price association varies steeply with PBU size: 0.214 (Q1),
0.098 (Q2), 0.045 (Q3), 0.017 (Q4)---a 12.5$\times$ gradient
consistent with Proposition~\ref{prop:optimal_m}
($\partial m^* / \partial \theta_k < 0$). The gradient motivates
CF3 (Section~\ref{sec:counterfactual}): concentrating oversight on
small PBUs maximizes the deterrence return.

\begin{figure}[htbp]
\centering
\includegraphics[width=0.85\textwidth]{fig_07_threshold_stability.pdf}
\caption{FL--price association across IQR multiplier thresholds.}
\label{fig:threshold_stability}
\end{figure}
